% This document was generated by an LLM.
% Content fragment; compile intro-analysis-self-check.tex.
% Coverage: Hunter (2012), Chapter 2, printed pp. 11--19.
% Clarifications of source slips:
% - A punctured neighborhood excludes its center; mere containment of a
%   deleted interval, as in the first sentence on p. 12, is insufficient.
% - One-sided limits require accumulation from the relevant side to avoid
%   vacuous conditions (Definition 2.9).
% - In the converse proof of Theorem 2.4, negate convergence to the given L.
% - In the proof of Theorem 2.23, use |x-c|, and L-M-2epsilon = 0.

\section{Limits of Functions}

\subsection{Limits}
\textit{Source: \S2.1, pp.~11--14.}

\begin{enumerate}
  % Definition 2.1.
  \item State the \(\epsilon\)-\(\delta\) definition of
    \(\lim_{x\to c}f(x)=L\) for \(f:A\to\R\).
    Include the hypothesis on \(c\), the domain restriction, and the
    order of all quantifiers.

  % Discussion following Definition 2.1.
  \item Must \(c\) belong to the domain for \(\lim_{x\to c}f(x)\) to
    exist? If \(f(c)\) is defined, does changing that value change the
    limit? Explain using the definition.

  % Definition 2.1; needed explicitly in the proof of Theorem 2.4.
  \item Write the precise negation of
    \(\lim_{x\to c}f(x)=L\), for a specified \(L\in\R\).
    How does this differ from saying that \(f\) has no finite limit at \(c\)?

  % Observation following Definition 2.1.
  \item Express convergence of \(f(x)\) to \(L\) in terms of a limit
    of \(|f(x)-L|\), and explain the equivalence directly from the definitions.

  % Example 2.2.
  \item For \(A=[0,\infty)\setminus\{9\}\), find
    \[
      \lim_{x\to9}\frac{x-9}{\sqrt{x}-3}.
    \]
    Prove your answer from the definition by giving an explicit
    \(\delta\) for each \(\epsilon>0\).

  % Paragraph preceding Definition 2.3, with the center excluded explicitly.
  \item Define a punctured (or deleted) neighborhood of \(c\) and a
    punctured \(\delta\)-neighborhood of \(c\). Must such a neighborhood
    contain \(c\)? Must it be an open set?

  % Definition 2.3 and its explanation.
  \item State the neighborhood formulation of \(\lim_{x\to c}f(x)=L\).
    Explain its equivalence with the \(\epsilon\)-\(\delta\) formulation
    in both directions.

  % Theorem 2.4.
  \item State the sequential characterization of a function limit
    (Theorem 2.4). Which sequences must be considered, and why are terms
    equal to \(c\) excluded?

  % Forward direction of Theorem 2.4.
  \item Prove that \(\lim_{x\to c}f(x)=L\) implies the sequential
    condition in Theorem 2.4. Explain how the choices of \(\delta\)
    and then \(N\) fit together.

  % Converse direction of Theorem 2.4.
  \item Prove the converse in Theorem 2.4 by contraposition.
    Starting from failure of convergence to the proposed \(L\), construct
    a sequence \(x_n\to c\) for which \(f(x_n)\not\to L\).
    Why must the same positive error threshold work for every \(n\)?

  % Corollary 2.5.
  \item State the two sequential criteria in Corollary 2.5 for showing
    that a finite function limit does not exist.
    Is checking one sequence whose function values converge sufficient
    to prove that a function limit exists?

  % Example 2.6.
  \item Define the sign function \(\operatorname{sgn}\) on \(\R\).
    Use two sequences to show that it has no limit at zero.

  % Example 2.7.
  \item Show, using the sequential criterion, that \(1/x\) has no
    finite limit as \(x\to0\). Can a single sequence establish this?

  % Example 2.8.
  \item Show that \(\sin(1/x)\) has no limit as \(x\to0\) by
    constructing two explicit sequences tending to zero whose function
    values have different limits.
\end{enumerate}

\subsection{Left, right, and infinite limits}
\textit{Source: \S2.2, pp.~14--16.}

\begin{enumerate}[resume]
  % Definition 2.9; one-sided accumulation made explicit.
  \item State the definitions of the right-hand and left-hand limits
    of \(f:A\to\R\) at \(c\). What accumulation condition on \(A\)
    is needed for each side? What would go wrong if the relevant side
    contained no domain points sufficiently close to \(c\)?

  % Example 2.10.
  \item Find the two one-sided limits of \(\operatorname{sgn}x\) at
    zero and verify them from the definitions.

  % Definition 2.11.
  \item Define \(\lim_{x\to\infty}f(x)=L\) and
    \(\lim_{x\to-\infty}f(x)=L\). Include the necessary assumptions
    on the domain and the quantifiers.

  % Discussion following Definition 2.11.
  \item How can a limit as \(x\to\infty\) or \(x\to-\infty\)
    be rewritten using the substitution \(x=1/t\)? Specify the side
    from which \(t\) approaches zero and how the domain is transformed.

  % Example 2.12.
  \item Find and prove the limits of \(x/\sqrt{1+x^2}\) as
    \(x\to\infty\) and as \(x\to-\infty\).
    Why must the identity \(\sqrt{x^2}=|x|\) be handled carefully?

  % Definition 2.13.
  \item State precisely what it means for \(f(x)\) to diverge to
    \(+\infty\) or \(-\infty\) as \(x\to c\).
    Do these statements assert the existence of a real-valued limit?

  % Example 2.14.
  \item Prove directly from the appropriate definitions that
    \[
      \lim_{x\to0}\frac{1}{x^2}=+\infty,
      \qquad
      \lim_{x\to\infty}\frac{1}{x^2}=0.
    \]
    Give explicit choices of the required radius or cutoff.

  % Example 2.15; combined one-sided and infinite-limit definitions.
  \item Write the quantified definitions of
    \(\lim_{x\to0^+}f(x)=+\infty\) and
    \(\lim_{x\to0^-}f(x)=-\infty\).
    Apply them to \(f(x)=1/x\).

  % Example 2.15.
  \item Why does \(1/x\) tend to neither \(+\infty\) nor
    \(-\infty\) as \(x\to0\), even though it is unbounded in every
    punctured neighborhood of zero?

  % Example 2.16.
  \item Show that \(x^{-1}\sin(1/x)\) has neither a finite limit nor
    divergence to either sign of infinity as \(x\to0^+\) or \(x\to0^-\).
    Construct sequences that demonstrate its behavior. What follows for
    the two-sided limit?

  % Example 2.17.
  \item Define \(\lim_{x\to\infty}f(x)=-\infty\) and
    \(\lim_{x\to-\infty}f(x)=+\infty\) using only inequalities
    between real numbers. Prove these two statements for
    \(f(x)=1/x-x^3\).

  % Opening paragraph of Section 2.2.
  \item State sequential characterizations of a finite right-hand
    limit, a finite limit at \(+\infty\), and divergence to \(+\infty\)
    at a finite accumulation point. Specify the permitted sequences
    in each case.
\end{enumerate}

\subsection{Properties of limits}
\textit{Source: \S2.3, pp.~16--19.}

\subsubsection*{Uniqueness and boundedness}
\begin{enumerate}[resume]
  % Proposition 2.18 and its proof.
  \item State and prove uniqueness of a finite function limit
    (Proposition 2.18). Where is the accumulation-point hypothesis used?

  % Paragraph following Proposition 2.18.
  \item Give a domain \(A\) and a point \(c\) that is not an
    accumulation point of \(A\). Show that without the accumulation
    hypothesis, the \(\epsilon\)-\(\delta\) condition for a limit
    could hold for every \(L\in\R\).

  % Definition 2.19.
  \item Define boundedness of a function on \(B\subseteq A\) and
    boundedness on its domain. Restate the first definition using the
    image set \(f(B)\).

  % Example 2.20.
  \item Give a function that is unbounded on its domain but bounded on
    a smaller interval. For \(f(x)=1/x\) on \((0,1]\), find a bound
    on \([\delta,1]\), where \(0<\delta<1\).
    For \(g(x)=x^2\) on \(\R\), find a bound on a finite interval \([a,b]\).

  % Proposition 2.21 and its proof.
  \item State and prove the local boundedness consequence of a finite
    limit (Proposition 2.21). What bound follows by choosing
    \(\epsilon=1\)? Is boundedness on the entire domain asserted?

  % Contrasting Proposition 2.21 with Example 2.8.
  \item Does boundedness in a punctured neighborhood imply existence
    of a limit? Give a counterexample from this chapter.
\end{enumerate}

\subsubsection*{Algebraic properties}
\begin{enumerate}[resume]
  % Theorem 2.22.
  \item State the limit laws for scalar multiples, sums, products,
    and quotients (Theorem 2.22), including all hypotheses and the
    additional restriction in the quotient law.

  % Sum proof in Theorem 2.22.
  \item Prove the sum law directly from the \(\epsilon\)-\(\delta\)
    definition. How do you choose a radius that works for both functions?

  % Product proof in Theorem 2.22.
  \item Prove the product law directly from the definition.
    Explain how to split \(f(x)g(x)-LM\) into two error terms and
    where local boundedness is used. Ensure that your choices also
    work when \(L=0\).

  % Quotient law in Theorem 2.22; remaining proofs are exercises in the notes.
  \item If \(g(x)\to M\ne0\), show that \(|g(x)|\) is bounded
    away from zero on the domain sufficiently close to \(c\), with
    \(x\ne c\). How does this ensure that the quotient is defined
    there? Use an estimate for \(|1/g(x)-1/M|\) to justify the quotient law.

  % Introductory paragraph of Section 2.3 and proof of Theorem 2.22.
  \item Explain how the sequential characterization of function limits
    transfers the algebraic limit laws for sequences to functions.
\end{enumerate}

\subsubsection*{Order properties}
\begin{enumerate}[resume]
  % Theorem 2.23 and its proof.
  \item State and prove the order-preservation theorem for finite
    limits (Theorem 2.23). In the contradiction proof, how do you
    choose a positive error tolerance from the two proposed limits?

  % Non-strict inequalities emphasized at the start of Section 2.3.3.
  \item If \(f(x)<g(x)\) for every domain point near \(c\), must
    their limits satisfy a strict inequality? Give a counterexample.

  % Theorem 2.24.
  \item State the squeeze theorem (Theorem 2.24), including the
    assumptions that guarantee existence of the middle function's limit.

  % Standard application following Theorem 2.24.
  \item Suppose \(|f(x)|\leq g(x)\) near \(c\), with \(x\ne c\),
    and \(g(x)\to0\). What follows about \(f(x)\), and how does
    the squeeze theorem give it?

  % Example 2.25.
  \item Give an example showing why convergence of the two bounding
    functions to the same limit is essential in the squeeze theorem.
    Can boundedness of the middle function alone replace this hypothesis?
\end{enumerate}
